\documentclass[11pt]{article}
\usepackage[utf8]{inputenc}
\usepackage[T1]{fontenc}
\usepackage{lmodern}
\usepackage{geometry}
\usepackage{hyperref}
\usepackage{longtable}
\usepackage{booktabs}
\usepackage{array}
\usepackage{xcolor}
\usepackage{colortbl}
\usepackage{enumitem}
\usepackage{amsmath,amssymb}
\usepackage{textcomp}
\usepackage[htt]{hyphenat}
\usepackage{titlesec}
\titlespacing*{\paragraph}{0pt}{0.5em}{0.6em}
\sloppy

\geometry{margin=2.5cm}

\definecolor{headerblue}{HTML}{1a1a2e}
\definecolor{rowgray}{HTML}{f6f8f8}

\hypersetup{
  colorlinks=true, linkcolor=headerblue, urlcolor=headerblue,
  pdftitle={Grassroots Networking: NAT Detection and Hole-Punching Survey},
  pdfauthor={Dan Bachar}
}

\newcolumntype{L}[1]{>{\raggedright\arraybackslash}p{#1}}
\newcommand{\code}[1]{\texttt{\small #1}}

\setlist[itemize]{leftmargin=1.4em, itemsep=0.2em, topsep=0.3em}
\setlist[enumerate]{leftmargin=1.6em, itemsep=0.2em, topsep=0.3em}
\setlength{\parskip}{0.4em}
\setlength{\parindent}{0pt}

\title{\vspace{-1.2cm}Grassroots Networking\\[0.2em]\large NAT Detection and Hole-Punching Survey}
\author{Dan Bachar \\ \small\texttt{dan.bachar@campus.technion.ac.il}}
\date{}

\begin{document}
\maketitle
\vspace{-0.8cm}

\section{Introduction}
\label{sec:intro}

Grassroots Networking (GN) is a peer-to-peer messaging transport that moves packets between mobile
devices over two mediums: Bluetooth Low Energy (BLE) for peers in physical proximity, and UDP over the
Internet for remote peers. The Internet relation assumes successful NAT traversal across heterogeneous
home, mobile and enterprise networks, and the decision space for that relation,
whether to dial directly with hole punching, to coordinate through a rendezvous point, or to fall
back on a relay, depends on how NAT is actually deployed in the networks our users inhabit. We have
never measured any of this: every constant in the transport, from the two-second punch to the
ten-second connection timeout, was chosen by intuition and not derived from data.

Public NAT measurements exist, but they are dated and skewed towards web and CDN workloads observed
from desktop and server vantage points~\cite{Jennings2005,Guha2005}. Mobile networks are the least
covered by those measurements, and at the same time the least favorable to traversal, as a handset
behind a carrier-grade NAT (CGNAT) cannot observe its own outer
address, cannot accept unsolicited inbound traffic, and may be assigned a different external port for
every destination it contacts, which defeats hole punching entirely. Desktop tooling cannot observe
any of this, as it does not sit behind the carrier; measuring it requires code running on handsets
carrying commercial SIMs, over the carrier networks those SIMs traverse.

In addition to the measurement gap, a second problem lies in the architecture of reachability
detection itself. The established way for a peer to learn whether it is reachable
is to ask other nodes to dial it back, as libp2p's autoNAT does~\cite{autonat}. Although autoNAT is
nominally decentralized, in practice it depends on a small set of well-known dial-back nodes, which
reintroduces the privileged infrastructure that a grassroots system is designed to do without.
GN already maintains a social graph of mutually authenticated friends, and we therefore ask whether
that graph can serve as the witness set instead, and what accuracy we lose by replacing designated
servers with friends.

In this project we design and validate a NAT detection algorithm that uses social-graph peers as
\emph{dial-back witnesses}, evaluate it against a reference implementation on a testbed under our control,
and use the resulting measurements to derive a decision rule for the transport. The results determine
whether rendezvous anchors are needed by the median user or only as a fallback, whether to invest
further in the UDP relation or push BLE harder, whether a relay is ever necessary, and how strongly to
prefer IPv6 over IPv4.

\section{The current transport and its blind spots}
\label{sec:current}

We first describe the IP path as it is implemented today, since the gaps in
Section~\ref{sec:gaps} follow directly from it. References are to files in the
\code{grassroots\_networking} repository.

The transport binds one wildcard UDP socket per address family, both on port 0, so the port is
assigned by the operating system and changes on every rebind
(\code{udp\_transport\_service.dart:1009}). A UDX multiplexer wraps each socket and owns all reads
from it. A device learns its own address in two ways: by an HTTPS request to an external reflection
service (\code{public\_address\_discovery.dart:40}), and through \code{ADDR\_REFLECT}, a signaling
message in which a well-connected peer reports the source address it observed on an incoming session
(\code{signaling\_service.dart:760}). No STUN client exists in the repository.

The address a device advertises is the externally discovered public IP paired with its own locally
bound port (\code{public\_address\_discovery.dart:97}). This reveals the public IP but never the
NAT-mapped port, and the only source of the mapped port, \code{ADDR\_REFLECT}, requires a session that
already works. Self-classification is a single boolean, \code{isWellConnected}, computed as a
routability test over that advertised string. The routability test itself is correct and excludes
100.64/10 among other reserved ranges (\code{address\_utils.dart:138}), but a handset behind a CGNAT
never observes its inner address, so the exclusion never applies to it.

Peers advertise address candidates in a signed \code{ANNOUNCE} that repeats every ten seconds; the
same timer drives liveness, reconnection and the queued-message drain, and there is no dedicated NAT
keepalive. Candidate selection scans every local-remote pair and returns exactly one, ranked
link-local, then IPv6, then IPv4 (\code{connection\_service.dart:48}); nothing probes candidates in
parallel, validates them, or falls back to the second-best pair after a failed attempt. The hole punch
is a 36-byte packet carrying a magic value and the sender's public key, sent for two seconds at
200\,ms intervals to a single address and port, without port prediction
(\code{hole\_punch\_service.dart:117}). A closed-loop variant exists but has no production caller and
cannot function as wired, since it reads the raw socket that the multiplexer already owns. The dial
itself is a single UDX attempt with a ten-second timeout and no retry inside the transport.

\subsection{What we cannot currently determine}
\label{sec:gaps}

\textbf{No NAT classification}: no code determines mapping or filtering behavior, and the boolean
that stands in for it is systematically wrong for the devices that matter most. A handset behind a
CGNAT that can reach the reflection service advertises a globally routable address, and is therefore
classified as well-connected both by itself and by its peers. Three behaviors follow from this false
positive: the pre-connect punch is suppressed for peers that appear public
(\code{grassroots\_network.dart:3369}), and is thus skipped for precisely the peers that require it;
the device begins acting as an address reflector and rendezvous facilitator for its friends; and the
user interface reports that the device is well-connected. Nothing in the system distinguishes having a
public address from being reachable at it.

\textbf{No mapped-port discovery}: the port we instruct peers to dial is a guess whenever no session
has yet succeeded, which is the case that matters for a first contact.

\textbf{No use of reflector disagreement}: several rendezvous servers and well-connected friends may
each reflect the same socket, which is precisely the raw material of an RFC 5780 mapping test, but the
handler overwrites the stored address and retains nothing, so two disagreeing reflections are
indistinguishable from a roaming event.

\textbf{No dial-back and no witnesses}: reachability from outside is asserted from an address prefix
and never tested, and the inbound-connection observation point that exists in the transport is never
read for a decision.

\textbf{No instrumentation}: punch outcomes are recorded as a four-value enumeration without
timestamps, and the reducer discards the failure reason, so neither latency nor robustness can be
computed after the fact. Hole punching is open-loop, and a refused mapping, a lost packet, a wrong
port and a slow handshake are indistinguishable in the logs.

A concrete instance illustrates the combined effect. Consider two friends who have previously paired
over BLE and are now remote from one another: one on a handset behind a carrier CGNAT with an
IPv4-only path, the other behind a consumer router at home with IPv4 and a working IPv6 prefix. The
first device advertises the carrier's outer address paired with its own local port, and both devices
compute that it is well-connected, which is wrong both in the port and in the reachability it implies.
Candidate selection then prefers the IPv6 candidate of the second device, because an IPv6 socket is
bound on the first even though it has no IPv6 transit, and the single attempt times out after ten
seconds without falling back to the IPv4 pair. The pre-connect punch is suppressed on both sides
because each observes the other as public, and had it fired it would have targeted the local ports.
Nothing is recorded, and the failure is indistinguishable from the peer being offline. When the same
code succeeds, on a port-preserving NAT where the guessed port happens to be correct, the system
learns nothing from that either.

\section{Research questions}
\label{sec:rqs}

\textbf{RQ1}: can social-graph peers replace designated dial-back servers? We specify a witness-based
algorithm that recovers a peer's mapping and filtering behavior, and quantify its agreement with an
RFC 5780 reference implementation.

\textbf{RQ2}: what does decentralization cost? We measure how the false-negative rate and the
time-to-verdict vary with the witness count $k$ and the byzantine fraction $f$, for both colluding and
independent adversaries, and determine whether the gap to a designated-server baseline closes as $k$
grows. We also evaluate the availability constraint, $P(\geq k$ witnesses online and dialable$)$,
which we expect to dominate the accuracy question in practice.

\textbf{RQ3}: which NAT behaviors occur in the networks our users inhabit, and how prevalent are
they? We measure the tuple of Section~\ref{sec:measurand} per stratum on our own devices, and estimate
population prevalence by weighting strata with published subscriber counts, as described in
Section~\ref{sec:prevalence}.

\textbf{RQ4}: what should the transport do at dial time? We derive a decision rule over the inputs
the transport holds at the moment it must decide, namely to attempt directly, to request mediation,
or to abort, and validate it by replay.

\subsection{Measurand}
\label{sec:measurand}

Instead of a NAT type, and following RFC 4787 and RFC 5780~\cite{rfc4787,rfc5780}, we decided to
measure a tuple: \emph{mapping behavior}, which is endpoint-independent, address-dependent or
address-and-port-dependent; \emph{filtering behavior}, which is endpoint-independent (EIF),
address-dependent (ADF) or address-and-port-dependent (APDF); \emph{port preservation}; and
\emph{mapping lifetime}.
The orthogonal facts of CGNAT suspicion and IPv6 availability are recorded alongside. A single-label
NAT type may be derived from the tuple for presentation, but is never the measurand.

The tuple is measured and reported per address family and per translation path, which is either NAT44,
NAT64 with 464XLAT, or native IPv6. On a 464XLAT path the entity being characterized is the carrier's
translator, which the RFC 4787 vocabulary describes only partially, and on a native IPv6 path there is
typically no mapping at all and only stateful filtering, so mapping behavior and port preservation
are reported as not applicable rather than as absent. The translation path therefore splits strata and
is not an orthogonal attribute.

We note that autoNAT does not produce this measurand. Version 1 reports whether a node is publicly
reachable, and version 2 reports whether a given address is dialable, with a nonce in the dial-back
request~\cite{autonat}; neither emits a mapping or filtering label. The comparison against autoNAT is
therefore restricted to reachability, and the ground truth for classification is an RFC 5780 reference
server that we operate ourselves.

\section{Approach}
\label{sec:approach}

\subsection{Witness-based detection}

To characterize the path the transport actually uses, we decided to implement the probe inside
\code{grassroots\_networking}, under \code{lib/src/transport/}, and to send from the transport's
existing per-family socket instead of from a socket of its own, as a second socket receives a
different NAT mapping and thus characterizes a different path altogether. The UDX fork exposes a raw-packet callback for this purpose. Three
properties of the existing code constrain the implementation: the raw handler discards datagrams
shorter than 50 bytes and discards 36-byte punch packets, so probes must exceed 50 bytes and carry
their own magic value; the existing raw send path registers its targets in a map that is not cleared
on shutdown, so a separate send path is required to prevent the announce broadcast from reaching
measurement targets; and the closed-loop punch mode is inoperative for the read-ownership reason given
in Section~\ref{sec:current}.

Several components are reused instead of rebuilt: \code{hole\_punch\_service.dart} for both punch
modes; \code{address\_utils.dart} for its address-class labeler, which already excludes CGNAT, ULA,
Teredo and 6to4; \code{connection\_service.dart} for candidate scoring, on which the probe loop is
layered; \code{noise\_session\_manager.dart} for the authenticated channel; \code{peer\_link.dart}
together with the anchor's redemption handler as the template for a signed witness attestation; and
\code{bootstrap\_anchor/}, which is the only component that runs without Flutter and is therefore the
natural host for the reference server.

The two RFC 5780 dimensions are not equally reachable from social witnesses. Mapping behavior
requires the mapped address to be reflected from at least two witnesses at distinct addresses, plus
one at a distinct port on one of them; \code{ADDR\_REFLECT} already performs the reflection, and what
remains is to reflect from a second distinct address and to retain each observation as a
\code{(reflector, tuple, timestamp)} triple instead of overwriting it. Filtering behavior is the
harder half, as it requires an inbound packet from an endpoint the probe has never contacted, which no
single witness can produce for itself. We therefore specify a relayed dial-back, carried on the existing signaling plane inside Noise
sessions, where clear signaling is refused on both ends and the eleven types in use occupy
\code{0x05}--\code{0x0f}, so that the witness messages are purely additive from \code{0x10} upwards.
The probe sends witness $W_1$ a nonce, its own mapped address and the requested source class; $W_1$ instructs a second
witness $W_2$, which the probe has not contacted, to dial back; and $W_2$ dials the probe carrying the
nonce, which the probe verifies. Three variants of the dial-back source yield EIF, ADF and APDF
respectively: an already-contacted port on a contacted address, a new port on a contacted address,
and an entirely new address.

\subsection{Contamination}
\label{sec:contamination}

The central validity threat, which we refer to as \emph{contamination}, is that social witnesses are, by construction, the peers the device
already transmits to: an \code{ANNOUNCE} every ten seconds on every live connection, and roughly ten
punch datagrams per traversal attempt. A dial-back arriving from an endpoint that the probe has
previously contacted traverses a mapping the probe itself opened, and consequently measures the
probe's own behavior rather than the NAT's filtering rule, which biases results towards reporting
NATs as more permissive than they are.

We mitigate this by maintaining, for the duration of a test, a send history of every destination
address and port emitted on the measured socket, treating a witness source that appears in that
history as a logged discard condition, and requiring every dial-back in the filtering arm to originate
from an endpoint never contacted. We report the size of the effect by running the same test with and
without a prior punch towards the witness, since the resulting label difference quantifies the bias
that naive witness-based reachability tests incur.

\subsection{Aggregation and adversary model}

A dial-back carrying a client-chosen nonce cannot be forged, so a single honest witness is sufficient
to prove reachability, whereas a claimed failure can be fabricated or withheld at no cost. The
aggregation rule is therefore asymmetric: any verified success establishes reachability, and
unreachability is declared only after $k$ independent witnesses have failed. Instead, a symmetric majority vote would
discard that asymmetry and treat a fabricated failure as evidence.

We define the adversary by its byzantine fraction, whether its members collude, and whether it is
on-path or off-path, and we implement three defenses before the probe runs outside the bench. First,
the client-chosen nonce must be echoed in the dial-back, so that a witness cannot claim a dial it did
not perform. Second, an off-path dial must target the requester's observed source address, or else
carry a proof-of-cost comparable to that of autoNAT v2; we report what this costs in bytes and
latency, since it is both a metered-data consideration and a latency confound. Third, we apply
per-requester rate limiting and refuse to dial or punch any target that is not globally routable,
using the predicate that already exists in \code{address\_utils.dart} but is not currently applied on
the punch path. We additionally require two independent witnesses to agree before the stored public
address is modified. Without these defenses a witness protocol layered on the current code would be
weaker than autoNAT v1, whose anti-amplification stance was to refuse dialing any address other than
the requester's observed source.

All measurement traffic is confined to the testbed by a target allowlist enforced at the send path, so
that a witness-supplied address cannot cause a probe or punch towards a third party, and measurement
is never enabled in the production build.

\subsection{Prevalence estimation}
\label{sec:prevalence}

A convenience sample of our own handsets cannot be averaged into a prevalence figure directly, but it
does not need to be, because NAT behavior is primarily a property of the network path rather than of
the individual subscriber: every customer of one carrier, on the same access point name (APN) and PDN
type, traverses the same small set of gateway configurations. We therefore characterize each stratum
on the testbed and weight the strata by published population data, while accepting that with few
measurements per stratum the within-stratum variance is not estimable and can only be bounded by
assumption. This is a stratified census of strata and not a stratified sample, as there is no
probabilistic sampling within a stratum; consequently no confidence interval is computable, and every
interval we report is a sensitivity range over assumptions.

One subscriber-level determinant survives this argument and is handled explicitly: the handset stack
selects the requested PDN type and whether a translator is present, and since Android and iOS differ
here, two subscribers on one carrier and APN may traverse different translators. Each cellular stratum
is therefore measured on at least two handset stacks, one of each platform.

\textbf{Cellular strata}: one per carrier, APN and PDN type rather than one per carrier. These strata are the tractable half and
not the certain half: they are few and each is cheap to enumerate, but each carries a large weight,
which is precisely why an undetected error inside one is expensive, and why each is covered by at
least two SIMs and two handset stacks before its homogeneity is treated as established. Homogeneity within a carrier remains an assumption, and the
ways it fails are known: the APN tier, since a consumer prepaid SIM generally cannot reach corporate or
M2M gateways; the arrangement of a mobile virtual network operator, as a full MVNO operates its own
core and NAT and constitutes a separate stratum while a light MVNO does not; the regional gateway
pool; and the behavior of the address pool under exhaustion. Weights are taken from published
subscriber counts and market-share reports, cross-checked against per-AS end-user estimates.

\textbf{Residential strata}: one per ISP and per whether that ISP applies CGNAT to the subscriber,
with customer-premises equipment vendor and firmware family as a secondary split within the
non-CGNAT cells only. The ordering matters, because where CGNAT sits above the customer equipment the
behavior of that equipment is not observable from outside and the vendor axis explains nothing; it
also determines which weights are obtainable, as ISP subscriber share is published whereas equipment
share by firmware family is not. These strata are the weak half, since within the secondary split the
number of models, the drift of firmware and subscriber-configurable settings make homogeneity a
genuine assumption. Should per-ISP equipment weights prove unobtainable, we report residential results
as an interval across the observed cells, weighting only the ISP and CGNAT axis.

\textbf{Managed and public networks}: university, office and commercial hotspots, characterized as a
set and carried with a coarse and explicitly uncertain weight.

\textbf{Residual stratum}: tethered clients, whose NAT stacks above the carrier's; VPN and
private-relay users, who replace the measured path entirely and are over-represented in the
privacy-conscious population that GN attracts; fixed-wireless and satellite broadband, which is
residential in name but cellular in mechanism; and enterprise networks. We do not measure this
stratum, and carry it at its best available weight so that it enters the bound at both extremes. It is
what allows the weight table to sum to the denominator.

The denominator is stated once: consumer-APN mobile lines and residential broadband lines in the
countries we cover, at the date of measurement, and not devices, humans or GN users. Published
subscriber counts include machine-to-machine and data-only lines, which we net out where the source
permits and flag where it does not. Weights are taken from subscriber counts, cross-checked against
per-AS end-user estimates such as those published by APNIC, which estimate eyeball users per AS from
an ad-delivered sample and cannot be split into cellular and fixed for a carrier announcing both
under one ASN, and never from the size of the announced address space, which is uncorrelated with user population once CGNAT is present and would systematically
under-weight the strata of interest.

We report prevalence with its weight table, and accompany it with a sensitivity analysis in four arms:
under each independent published weighting; under a worst-case and best-case assignment of all
unmeasured and heterogeneity-uncertain mass to the extremes of the measurand domain, which is the only
arm that bounds the conclusion; under the assumption that unmeasured strata behave as the worst
observed stratum, which is a pessimistic variant and not a bound; and under perturbation of each
measured stratum's label across the range its homogeneity check observed. We name which envelope
dominates, and report the cellular and residential halves separately. Comparison against published NAT
and DCUtR measurements~\cite{dcutr} is performed only where the denominators reconcile, as those
studies estimate over libp2p nodes and traversal attempts and are skewed towards always-on desktop and
datacenter hosts.

\section{Experiments}
\label{sec:experiments}

The evaluation is conducted entirely on a testbed under our control, in two tiers. The first tier is a
synthetic bench whose labels are known by construction, and on which the correctness claim is
established. The second tier consists of our own handsets, carrying prepaid SIMs across the local
carriers, moved across the access networks we can physically reach. As the second tier carries no
labels, we report agreement with the reference implementation rather than accuracy, and treat
disagreements as a category with hypotheses such as mapping churn during a test, rate limiting at the
reference server, or multipath.

\subsection{Bench validation against a reference implementation}

\paragraph{Question.} Does the witness-based algorithm recover the mapping and filtering behavior of
a NAT whose configuration is known by construction, and how does its agreement compare with that of an
RFC 5780 reference server and of autoNAT?

\paragraph{Setup.} A reference server with two public addresses and two ports per family, verified to
advertise its alternate address and to honour change requests from it. The bench spans the mapping
axis with netfilter: the default masquerade rule reuses a mapping for a given source tuple and
presents as endpoint-independent, while \code{--random-fully} allocates a fresh port per flow and
presents as address-and-port-dependent. No netfilter option produces address-dependent mapping, so
that case is constructed explicitly with per-destination source NAT through nftables maps. We note
that \code{--persistent} controls the source address pool and does not belong to the mapping axis. The
bench additionally spans EIF, ADF and APDF filtering, includes two to three commercial routers, and
includes a double-NAT configuration emulating CGNAT.

\paragraph{Procedure.} Each configuration is measured by the reference server, by the witness-based
algorithm with three honest witnesses, and by autoNAT v1 and v2, back to back within one window. The
contamination guard of Section~\ref{sec:contamination} is then disabled and the filtering arm repeated,
with a punch issued towards the witness beforehand, to quantify the bias it prevents.

\paragraph{Measures.} Per-configuration labels from each method against the construction label;
wall-clock time to verdict; bytes transmitted and received per verdict, including the proof-of-cost
payload where applicable.

\paragraph{Output.} Confusion matrices against the construction labels for at least eight
configurations; agreement of the witness-based algorithm with the reference server; and the label
difference between contaminated and uncontaminated runs.

\subsection{Robustness to dishonest witnesses}

\paragraph{Question.} How do the false-negative rate and the time to verdict degrade as the witness
count and the fraction of dishonest witnesses vary, and does the asymmetric aggregation rule of
Section~\ref{sec:approach} behave as intended?

\paragraph{Setup.} Witness processes under our control on the bench, of which a fraction $f$ are
instructed to misbehave by claiming dial-backs they did not perform, by withholding successful ones,
and by colluding on a shared account of events. Since we operate the witnesses, the ground truth of
which witness lied is known by construction.

\paragraph{Procedure.} Sweep the witness count $k$ against $f \in \{0, 0.1, 0.3, 0.5\}$, for
independent and for colluding adversaries. The availability constraint is evaluated parametrically, by
sweeping the online fraction and the friend-count distribution, as the testbed cannot supply a
realistic distribution of its own.

\paragraph{Measures.} False-negative and false-positive rate per $(k, f)$; time to verdict; number of
dial-back requests issued per verdict.

\paragraph{Output.} Error rate against $k$ for each $f$, exhibiting the asymmetry of the aggregation
rule; time to verdict against $k$; and the parametric availability curve $P(\geq k)$.

\subsection{Testbed campaign across strata}

\paragraph{Question.} Which values of the measurand occur across the strata of
Section~\ref{sec:prevalence}, and how stable is each stratum under repetition and under variation of
the subscriber configuration?

\paragraph{Setup.} The handsets and SIMs of the testbed inventory, across the carriers and access
networks it names. Each cellular stratum is covered by at least two SIMs differing in plan or operator
arrangement, and by at least two handset stacks.

\paragraph{Procedure.} Each stratum is measured on at least two days, spanning a peak and an off-peak
window, with the screen on and off and the application in the foreground and the background. Roaming
and handover are triggered deliberately. Mapping lifetime is probed per network with the announce
timer quiesced for the duration of the probe.

\paragraph{Measures.} The full measurand per stratum, with the translation path; the pairwise
agreement between the subscriber configurations within a stratum, which serves as the homogeneity
check; interface type, radio technology, screen state and application state per record; and agreement
with the reference server.

\paragraph{Output.} A per-stratum table of the measurand; the homogeneity check per cellular stratum,
where a disagreement splits the stratum rather than being averaged; the list of strata that were not
reached, with their population weight, so that the coverage gap is quantified; and the weighted
prevalence estimate with the four-arm sensitivity analysis of Section~\ref{sec:prevalence}.

\section{Deliverables}
\label{sec:deliverables}

\begin{itemize}
  \item The detection algorithm: a specification with its adversary model, a reference implementation
        in \code{grassroots\_networking}, and its validation against the RFC 5780 implementation.
  \item A reachability comparison against autoNAT v1 and v2, measured on one device, network and
        window. The comparison is Android-only, as embedding the reference implementation on iOS is a
        project of its own, and we state this as a limitation.
  \item The testbed dataset: every bench configuration with its construction label, and every stratum
        with its measured tuple, released in full, as these are our own devices on our own SIMs and
        nothing must be withheld. We check the acceptable-use terms of each prepaid and residential
        plan beforehand, as consumer terms occasionally prohibit network probing.
  \item Prevalence estimates with the weight table, the denominator, the date of measurement and the
        sensitivity analysis, cellular reported separately from residential.
  \item A per-attempt instrumentation subsystem recording the attempt identifier, peer, candidate,
        address family, path taken, phase timestamps, outcome and failure kind, with a durable sink and
        the attempt identifier carried on the control messages so that both endpoints' records join.
  \item A decision rule for the transport, validated by replay on a held-out split.
  \item A technical report, targeting ANRW or PAM.
\end{itemize}

The primary claim is the detection algorithm, validated on labeled ground truth and exercised against
real access networks; the prevalence estimate, the reachability comparison and the decision rule are
supporting results. Broadening the testbed to further carriers and countries improves the estimate and
is pursued if time permits. Related work is written in the first month rather than the last, as a
longitudinal DCUtR campaign already covers a large number of traversal attempts~\cite{dcutr} and the
axes on which this work is defensible, namely cellular carriers and CGNAT observed from real handsets
and the decentralization of the witness set, must be positioned against it from the outset.

\section{Milestones}
\label{sec:milestones}

{\footnotesize
\begin{longtable}{L{1.5cm} L{6.2cm} L{7.3cm}}
\toprule
\rowcolor{headerblue}
\textcolor{white}{\textbf{Window}} & \textcolor{white}{\textbf{Work}} &
\textcolor{white}{\textbf{Exit criterion}} \\
\midrule
\endhead
Month 1 & Scope, measurand, adversary model, aggregation rule, related work, stratum design and
testbed inventory. & A design document containing the mapping from each test to the witness topology
it requires, the adversary model, the aggregation rule, the stratum definition with its weight sources
named, and an inventory of every handset, SIM, router and access network required, marking what is
available and what must be procured. \\
\rowcolor{rowgray}
Month 2 & Reference server on two addresses and two ports; the labeled bench, including the
address-dependent mapping case that netfilter cannot produce. & One command measures at least eight
configurations and emits a confusion matrix, with all three mapping classes present; agreement with
the construction labels of at least 95\%, each disagreement explained; the bench reproduces from a
clean checkout. \\
Month 3 & The detection algorithm: relayed dial-back, filtering test, contamination guard,
anti-amplification defenses, multi-witness aggregation. & At least 90\% agreement with the reference
server across the bench configurations using three honest witnesses; tests demonstrating that a
witness cannot flip a verdict by claiming a dial it did not perform, that a witness cannot induce a
dial towards a third party, and that a dial-back from a contacted endpoint is discarded; and the
measured effect of the contamination guard. \\
\rowcolor{rowgray}
Month 4 & The autoNAT baseline and the comparison harness; the sweep over witness count and byzantine
fraction. & One record per device, network and window carrying all verdicts with latency, bytes and
battery; error rate and time to verdict against $k$ for each $f$; and the parametric availability
curve. \\
Month 5 & The testbed campaign across all strata, including lifetime probes, roaming and handover. &
Every stratum characterized with its full tuple and its homogeneity check, each measured on at least
two days; unreached strata listed with their weight. \\
\rowcolor{rowgray}
Month 6 & Prevalence estimation with sensitivity analysis; the decision rule; the address-exposure
comparison; writing and submission. & Prevalence with weight table, denominator and date, the dominant
sensitivity envelope named, and no interval presented as a confidence interval; the decision rule as a
predicate table validated on a held-out split, with the time saved against the current fixed punch and
timeout; the number of distinct peers learning a device's address per day under the witness-based and
the server-based designs; and one script regenerating every figure from the released data. \\
\bottomrule
\end{longtable}
}

Procurement of addresses and firewall rules for the reference server begins in the first week, as it
gates the second month.

\section{Prerequisites and open decisions}

We expect basic mobile application development, basic networking, and familiarity with statistics and
a plotting library; experience with a peer-to-peer stack is an advantage. Dart and Flutter can be
learned during the first month, whereas the networking intuition is harder to acquire on the way.

Four decisions are taken before the work begins. First, whether the autoNAT baseline is the reference
implementation integrated as a sidecar, which is faithful but Android-only, or a reimplementation over
our own socket, which is comparable as a probe but makes the baseline our own protocol. Second,
whether the primary claim is the algorithm or the measurement, as both cannot be first-class within
six months. Third, who provisions the reference server, and whether its alternate addresses can be
drawn from the address block already reserved on the existing virtual machine. Fourth, whether to open an IPv4 rule on the
deployed rendezvous anchor, whose data-plane rule today admits \code{udp:9516} from \code{::/0} and
is therefore IPv6 only, which excludes the IPv4-only CGNAT population the study targets. Finally, the
testbed budget: each carrier is a candidate stratum, but each SIM covers only one combination of
carrier, APN and plan, and at least two such combinations and two handset stacks per carrier are
required before homogeneity can be checked instead of assumed.

We note separately that the unvalidated punch target, the absence of rate limiting on the signaling
path, the uncorroborated address reflection and the unbounded length field in the stream framer are
defects in the current implementation, independent of this project, and that their assignment should
be settled in advance.

\section*{Glossary}

\begin{description}[leftmargin=0pt, style=nextline, itemsep=0.2em]
  \item[464XLAT] A mechanism allowing IPv4-only applications to operate over an IPv6-only mobile
        network, by translating on the handset and again in the carrier network. Android and iOS
        differ in their support, which is why the handset stack can change what a measurement observes.
  \item[APN] The named gateway configuration a mobile device uses to reach the packet core. Prepaid,
        postpaid, corporate and machine-to-machine plans frequently use different ones, and therefore
        potentially different NAT behavior on the same carrier.
  \item[autoNAT] A libp2p service in which a node asks other nodes to dial it back in order to
        determine whether it is publicly reachable. Version 2 answers the question per address and
        includes a nonce. Both versions rely in practice on well-known dial-back servers, and neither
        classifies NAT behavior.
  \item[CGNAT] Network address translation applied by an ISP, typically a mobile carrier, sharing one
        public IPv4 address among many subscribers. It is common on cellular networks and hostile to
        peer-to-peer traffic, as even the public address is shared.
  \item[Dial-back] A request from peer A that peer B attempt a connection back to A. Success
        establishes that A is reachable from outside.
  \item[Filtering behavior] Which inbound traffic a NAT admits on an existing mapping:
        endpoint-independent, address-dependent, or address-and-port-dependent. Determining it
        requires a packet from an endpoint never contacted.
  \item[Hole punching] Two peers behind NATs transmit towards each other simultaneously, so that each
        NAT admits the reply to traffic it has already seen. It generally fails against
        address-and-port-dependent mapping.
  \item[Mapping behavior] How a NAT allocates the external port for outbound traffic: the same port
        for every destination, or a fresh one per destination, in which case no peer can predict where
        to send.
  \item[Mapping lifetime] How long a NAT retains an idle mapping, which determines the keepalive
        interval required to hold a path open.
  \item[MVNO] An operator reselling service on another operator's network. A full MVNO operates its
        own packet core and NAT and is a separate stratum; a light MVNO is not.
  \item[Port preservation] Whether a NAT retains the internal source port as the external one. Where
        it does, the address a device guesses from its own local port happens to be correct.
  \item[RFC 4787, RFC 5780] The former defines NAT behavioral terminology and requirements; the
        latter defines the tests that discover that behavior from outside, using a server with two
        addresses and two ports.
  \item[STUN] A protocol by which a peer asks a public server which address it appears to originate
        from. GN has no STUN client; \code{ADDR\_REFLECT} is its nearest equivalent.
  \item[UDX] The reliable multiplexed stream protocol over UDP used for the Internet relation,
        consumed as a fork published as \code{grassroots\_dart\_udx}.
\end{description}

\begin{thebibliography}{9}
\bibitem{Jennings2005} C. Jennings. NAT Classification Test Results. IETF Internet-Draft, 2005.
\bibitem{Guha2005} S. Guha and P. Francis. Characterization and Measurement of TCP Traversal through
  NATs and Firewalls. In \emph{Proc. ACM IMC}, 2005.
\bibitem{rfc4787} F. Audet and C. Jennings. Network Address Translation (NAT) Behavioral Requirements
  for Unicast UDP. RFC 4787, 2007.
\bibitem{rfc5780} D. MacDonald and B. Lowekamp. NAT Behavior Discovery Using Session Traversal
  Utilities for NAT (STUN). RFC 5780, 2010.
\bibitem{autonat} Protocol Labs. libp2p AutoNAT v1 and v2 specifications.
  \url{https://github.com/libp2p/specs/tree/master/autonat}.
\bibitem{dcutr} D. Seemueller et al. A Longitudinal Measurement of Hole Punching in the IPFS Network.
  arXiv:2510.27500, 2025.
\end{thebibliography}

\end{document}
